\documentclass[10pt,a4paper]{amsart}
\usepackage[margin=19mm]{geometry}
\usepackage{amsmath,amssymb,mathtools}
\usepackage[scr=rsfs,cal=euler]{mathalfa}
\usepackage{xfrac}
\usepackage[unicode,hidelinks,bookmarks=false]{hyperref}
\newcommand{\ie}{i.e.\ }
\newcommand{\cf}{cf.\ }
\newcommand{\resp}{respectively\ }
\newcommand{\Subsection}{\subsection}
\providecommand{\nobreakdash}{\nobreak\hbox{-}}
\setlength{\emergencystretch}{2em}
\multlinegap0pt
\usepackage{bm}
\usepackage{bbm}
\usepackage{dsfont}
\usepackage{xspace}
\usepackage{cancel}
\usepackage{mathtools}
\usepackage{amsthm}
\makeatletter
\@ifundefined{theorem}{\newtheorem{theorem}{Theorem}}{}
\makeatother
\newcommand{\G}{X}
\title[Edge-transitive token graphs as covers]{Edge-transitive token graphs as covers}
\author{Sergio G. G\'omez-Galicia, Octavio B. Zapata-Fonseca}
\date{Source: arXiv:2505.20632v1 (CC BY 4.0)}
\begin{document}
\maketitle

\noindent\emph{Source and licence.} Edge-transitive token graphs as covers. Version arXiv:2505.20632v1 (CC BY 4.0), distributed under the Creative Commons Attribution 4.0 International licence, as recorded in the arXiv licence metadata for this exact version and shown on the abstract page https://arxiv.org/abs/2505.20632v1. Full rights notice: licenses/arxiv-2505.20632-CC-BY-4.0.txt.\par\smallskip

\noindent\textit{Editorial scope.} Counted unit: the Theorem whose source label is \texttt{theo\_1} in the article (source0.tex lines 136--139, source label \texttt{theo\_1}), delivered together with the article's own complete original proof of it (source0.tex lines 140--189, one \texttt{proof} environment of 50 lines). No mathematics is written, repaired or re-derived: statement and proof are the article's own text line for line. Required context retained uncounted: none. Every definition, display and label of those lines is retained unchanged. Omitted from this package and not redistributed: 1--135, 190--298. Nothing inside a retained excerpt is omitted or altered: every retained body is delivered exactly as the article's lines are written, so no omission receipt of any kind accompanies this package. Citations: the retained excerpts carry no citation command at all, so no bibliography entry of the article is transcribed and no reference is redistributed. Reference adaptation: none was needed and none was made; every reference inside the delivered unit resolves inside it, so no unresolved reference or citation is produced. Printed slips kept literal: no printed word, symbol, hypothesis or number of the counted statement or of its proof was altered. Layout and class: the article's own class is \texttt{article} with packages including amsmath, amssymb, amsthm, amsfonts, babel, comment, enumerate, hyphenat, graphicx, hyperref; the wrapper is the lane's pinned \texttt{amsart} editorial wrapper at 10pt on A4, which restates the theorem-like environments the retained text needs and adds the article's own macro definitions that the retained text needs (\texttt{\textbackslash G}), with extra wrapper packages (bm, bbm, dsfont, xspace, cancel, mathtools). The delivered numbering is the unit's own. Assets: no retained span contains a figure environment, a graphics-inclusion command, a PDF-page inclusion, a table or a TikZ picture; no image file is copied into this package, the package declares no asset dependency and no image file is redistributed with it. Licence: version arXiv:2505.20632v1 of this article is distributed under the Creative Commons Attribution 4.0 International licence, as recorded in the arXiv licence metadata for this exact version and shown on the abstract page https://arxiv.org/abs/2505.20632v1. A full rights notice is delivered at licenses/arxiv-2505.20632-CC-BY-4.0.txt. Characters: every retained excerpt is pure ASCII, so the delivered engine, XeLaTeX with its default Latin Modern text font, reports no missing character.
\begin{theorem} \label{theo_1}
Let $n$ be an even number. 
The token graph $F_2(K_n)$ is isomorphic to the covering graph $\G^{(\alpha,\omega)}$ of a combined base graph $\G$ on $n/2$ vertices and with combined assignment on the group $\mathbb{Z}_{n}$.
\end{theorem}

\begin{proof}
We construct the combined based graph $\G$ on the group $\mathbb{Z}_{n}$ as follows. 
Since $n$ is even, $n/2$ is an integer. 
Let the $n/2$ vertices of $X$ be enumerated as
\[
x_1,x_2,\dots,x_{n/2}.
\]
The map $\omega$ that sends each vertex $x_i$ to a subgroup of $\mathbb{Z}_n$  is defined by 
\[
	\omega(x_i):=
	\begin{cases}
        0\mathbb{Z}_n=\{0\} & \textnormal{ if }  1\leq i < n/2, \\
		i\mathbb{Z}_n =\{0,i\} & \textnormal{ if }  i=n/2.
	\end{cases}
  \]  
 We remark that $\G$ contains multiple edges and loops. Now we define the voltage assignment $\alpha$ in $\G$: 
\begin{itemize}
    \item If $1\leq i,j\leq n/2$, $i\neq j$, then there is an edge $e_{i,j}:=x_ix_j$ such that $\alpha(e_{i,j})=0$.
    \item If $1\leq i < n/2$, then $x_i$ has a loop $x_ix_i$.
    \item If $1\leq i < n/2$ and $j>i$, then
    \begin{itemize}
        \item Each vertex $x_i$ has a loop $x_ix_i$ such that $\alpha(x_ix_i)=1$.
        \item There are the following edges: $e'_{i,j}=x_ix_j$, $f_{i,j}=x_ix_j$ and $f'_{i,j}=x_ix_j$, such that $\alpha(e'_{i,j})=i$, $\alpha(f_{i,j})=n-j+i$ and $\alpha (f'_{i,j})=n-j$.
    \end{itemize}
\end{itemize}
 
 Now we show that the covering graph $\G^{(\alpha,\omega)}$ is isomorphic to $F_{2}(K_n)$. 
 First observe that 
 \[|V^{(\alpha,\omega)}|=|V(F_2(K_n))|,\]
 since the number of vertices of $\G^{(\alpha,\omega)}$ is given by
\[
\sum_{i=1}^{n/2}[\mathbb{Z}_n:\omega(x_i)]= \sum_{i=1}^{n/2-1}[\mathbb{Z}_n:\omega(x_i)]+ [\mathbb{Z}_n:\omega(x_{n/2})]=\left(\frac{n}{2}-1\right)n+\frac{n}{2}=\binom{n}{2}.
\]

Now we define the map $\varphi$ from $\G^{(\alpha,\omega)}$ to $F_{2}({K_n})$ as follows.
\begin{itemize}
    \item If $ 1\leq i< n/2$, then
\[
\varphi\left((x_i,\{j\})\in V^{(\alpha,\omega)}\right)= \{1+j\mod{n},1+j+i \mod{n}\}%\in V(F_2(K_n)).%\quad \text{ for all } 0\leq j\leq n-1.
\]
for all $0\leq j\leq n-1$.
    \item If $i=n/2$, then 
 \[
\varphi\left((x_{i},\{j,n/2+j\}\in V^{(\alpha,\omega)}\right)= \{1+j \mod{n},1+j+n/2 \mod{n}\}%\in V(F_2(K_n)). % \quad \text{ for all } 0\leq j\leq n-1.
 \]
 for all  $0\leq j\leq n-1$.
\end{itemize}

 By construction of the combined voltage assignment $(\alpha, \omega)$ on the base graph $\G$, it can be seen that the map $\varphi$ defined above is an isomorphism between the covering graph $\G^{(\alpha,\omega)}$ and the token graph $F_2(K_n)$.
\end{proof}

\end{document}
